\begin{result}{F.3.30}{An exact invariant reduction for the higher-genus single-line problem}{Proved parametrization; mixed-matrix conclusion retained open}
Fix $g\geq4$, write $V=\langle a,b\rangle\perp W$ with
$\langle a,b\rangle=1$, and put $s=g-1\geq3$.
Let $S=\operatorname{Sp}(W)\times\mathbb G_m\subset
\operatorname{Sp}(V)$, where the torus sends $a\mapsto ta$,
$b\mapsto t^{-1}b$ and fixes $W$.
The actual universal image $U_{g,6}$ and the tangent space
$P_a=\{D:D(a)\in[\mathbb L_6,a]\}$ are $S$-stable.
The quotient $q_{g,6}$ is trivial as an $S$-representation.

The actual invariant projector, the215 decorated Dynkin input
span and five pure-complement inputs, the coupled scalar-polynomial
mixed equations, and the five-detector inclusion below are defined
and proved for every such genus. The mixed matrix and its required
contraction-image conclusion remain uncomputed.
\end{result}
\begin{proof}
We construct each map on the actual representation and keep all mixed equations coupled.
\paragraph{1. The actual invariant projector.}
Finite-dimensional complete reducibility supplies the projector
$e_S$ to the trivial $S$-isotypic component. Equivalently, it
lies in the finite algebra generated by the actual rational
representation matrices of $S(\Q_\ell)$, and is zero on every
nontrivial simple summand. It therefore preserves every
$S$-submodule, in particular
\[
 e_S(U_{g,6}\cap P_a)\subset (U_{g,6}\cap P_a)^S,
                 \qquad q_{g,6}e_S=q_{g,6}.                 \tag{LI.1}
\]
This is a linear projector on the actual universal image;
no individual Galois element is asserted to realize it.
It also preserves the full mixed-normalizer equations because
those equations hold on all of $U_{g,6}$.

\paragraph{2. Cyclic reconstruction of the positive constrained blocks.}
Identify a special degree-six derivation with its tensor
\[
 T(D)=\sum_i(x_i\otimes D(y_i)-y_i\otimes D(x_i))
                       \in V\otimes\mathbb L_7(V).
\]
Writing $\rho$ for the cyclic rotation moving the first letter
to the end, specialness is exactly $(1-\rho)T=0$.
The $S$-invariant tensor space decomposes by the number $m$
of occurrences of each of $a,b$, with $0\leq m\leq4$.
This is a decomposition of the ambient invariant tangent
space; its projections are not claimed to preserve $U_{g,6}$.

For $m\geq1$, the tangent equation is $D(a)=[A,a]$, with
$A\in\mathbb L_6(V)$ of multiplicities $m-1,m-1$ in $a,b$
and $8-2m$ letters from $W$. The associative centralizer
$\Q_\ell[a]$ and its primitive part $\Q_\ell a$ show that
$\operatorname{ad}(a)$ is injective in this degree. Thus $A$
is uniquely determined. Since $D$ is $S$-invariant, the same
injectivity proves that $A$ is $S$-invariant.

Let $\pi_b$ retain exactly the terms beginning in $b$, and
let $\mathcal S_8=\sum_{j=0}^7\rho^j$. Every cyclic tensor
in this multidegree satisfies
\[
                  \mathcal S_8\pi_b T=mT.
\]
This counts the $m$ occurrences of $b$ and remains valid
for periodic words. Since $\pi_bT=-b\otimes[A,a]$, the
exact reconstruction map is
\[
 \Psi_m(A)=-\frac1m\mathcal S_8(b\otimes[A,a]).             \tag{LI.2}
\]
The required parameter subspace is cut out by the two
literal linear equations
\[
 (1\otimes\operatorname{Dyn}_7)\Psi_m(A)=7\Psi_m(A),
          \qquad \pi_b\Psi_m(A)=-b\otimes[A,a].            \tag{LI.3}
\]
The first says that the suffix is in the free Lie algebra;
the tensor is already cyclic, so it is special. The second
gives the actual tangent equation. Therefore (LI.2) maps
the kernel of (LI.3) bijectively onto the entire $m$-block
of the invariant tangent space. No surface-quotient map
is used in this parametrization.

\paragraph{3. A finite genus-independent input span.}
The invariant tensors before imposing the Lie condition on
$A$ have $m-1$ copies each of $a,b$ and $4-m$ copies of the
symplectic pairing tensor $\omega_W$. Choose the slots for
the fixed letters and pair the remaining slots. The four
numbers of decorated six-slot diagrams are
\[
 \begin{array}{c|r|r|r}
 m&\#a=\#b&\#\omega_W&\#\text{diagrams}\\\hline
 1&0&3&15\\
 2&1&2&90\\
 3&2&1&90\\
 4&3&0&20.
 \end{array}                                                \tag{LI.4}
\]
Applying the degree-six Dynkin map to these $215$ tensors
spans all possible $A$, because its restriction to
$\mathbb L_6$ is multiplication by six. The invariant
tensor theorem gives precisely these diagram spans
\cite[Proposition~2.1, p.~8]{AuditMoritaSakasaiSuzuki2015v2}. They
are independent before the Dynkin map whenever $s\geq3$,
because there are at most three $W$-pairs. The target
diagrams in (LI.3) also have at most three $W$-pairs.
All permutations, signs, Dynkin operations and consistency
equations are consequently the same finite integer
diagram equations for every $s\geq3$.

For $m=0$, the tensor is entirely in $W$ and the space is
$\mathfrak h_{s,6}^{\operatorname{Sp}(W)}$. It has dimension
five for every $s\geq3$, with the explicit coherent
chord-diagram basis $v_1,\ldots,v_5$ in MSS2015,
printed pp.~25--27\cite{AuditMoritaSakasaiSuzuki2015v2}. Thus the entire invariant tangent
input is parametrized by the five pure-$W$ vectors and
the constrained Dynkin images of the $215$ diagrams.
Relations among those $215$ images are retained; this
statement is a span and an exact kernel construction,
not a declaration of $220$ independent coordinates.

\paragraph{4. Mixed equations stay coupled across the blocks.}
On the sum $D=D_0+\cdots+D_4$, form the actual linear map
\[
 B(D):\Lambda^3V\longrightarrow
       (V^{\otimes7})_{\rm cyc}\oplus V,\qquad
 \delta\longmapsto
   \bigl(\operatorname{ES}_7([D,\delta]),R([D,\delta])\bigr).
                                                               \tag{LI.5}
\]
Every actual $D\in U_{g,6}$ lies in its kernel, using the
actual geometric-normalizer maps of F.3.24 and the opened
primary formula for $R$\cite[Proposition~4.1, p.~11]{AuditMoritaSakasaiSuzuki2026v1}. The equations
are imposed on the sum. No individual $D_m$ is declared
to be an actual universal arithmetic vector.

The curried mixed map is an invariant tensor for $S$.
Pair it with invariant dual tensors, built from $a,b$
and symplectic contractions on $W$. Each resulting scalar
row is a necessary equation in every genus. These scalar
rows also separate the actual invariant target: the
invariant isotypic projector is compatible with tensor
duality, and the invariant tensor theorem spans the dual
invariants. This formulation accounts for low-genus
relations in output diagrams; formal coefficient
independence of an arbitrary ten-slot output is not
assumed at $s=3$ or $4$.

Each closed $W$-loop in a scalar contraction supplies a
factor $\dim W=2s$, with the orientation sign determined
by the symplectic form. Thus every specified scalar
row has an exact polynomial in $s$ obtained by the
finite contraction diagram. The input construction
does not by itself determine the rank or the kernel
of these mixed equations.

\paragraph{5. Five exact output coordinates suffice.}
MSS2015 Proposition~7.2, printed pp.~26--27\cite{AuditMoritaSakasaiSuzuki2015v2}, identifies
the five complete-contraction detector $K_{\mathcal D}$
with an isomorphism on
$\mathfrak h_{g,6}^{\operatorname{Sp}(V)}$ for all $g\geq3$.
Its matching indices among the $105$ diagrams are
$1,2,3,8,14$. Our raw pairing products differ from the
source's signed $\alpha_C$ coordinates by the explicit map
$\operatorname{diag}(1,-1,1,-1,-1)$. Indeed, the sign for
$C=(i_1j_1)\cdots(i_4j_4)$ is the permutation sign of
$(i_1,j_1,\ldots,i_4,j_4)$, exactly as defined on p.~7
of the source. All span and injectivity assertions are
transported through this invertible diagonal map. Let $Z_1,Z_2$ be the actual geometric
brackets of F.3.29, embedded using three chosen handles.
Their contraction vectors remain unchanged under the
symplectic inclusion, and their two-coordinate minor
is $-10752$.

The full symplectic invariant projector $e_G$ preserves
$U_{g,6}$ and $J_6$, and $q_{g,6}e_G=q_{g,6}$.
It is used here only on the output argument and is not
claimed to preserve tangency. Since $K_{\mathcal D}e_G
=K_{\mathcal D}$ and the detector is injective on full
invariants, the following unconditional linear inclusion
holds on the actual universal image:
\[
 \ker\!\left(U_{g,6}\xrightarrow{K_{\mathcal D}}
  \Q_\ell^5/\operatorname{span}\{K_{\mathcal D}(Z_1),
                                      K_{\mathcal D}(Z_2)\}\right)
                 \subset\ker q_{g,6}.                    \tag{LI.6}
\]
Indeed, the relevant equality of five coordinates gives
an equality after $e_G$ with a linear combination of
$e_GZ_1,e_GZ_2\in J_6$. Applying $q_{g,6}$ proves (LI.6).
This is an invariant map on the original genus-$g$
derivation space and makes no projection to a smaller
arithmetic surface quotient.

\paragraph{Current exact boundary.}
The construction (LI.1)--(LI.6) is proved. The scalar
mixed matrix on its coupled invariant input has not
yet been evaluated and its required contraction-image
inclusion is not established. No all-genus single-line
$q_6$ theorem or full LM6 A.22 resolution is asserted.

\end{proof}
